\documentclass[10pt,a4paper]{amsart}
\usepackage[margin=19mm]{geometry}
\usepackage{amsmath,amssymb,mathtools}
\usepackage[scr=rsfs,cal=euler]{mathalfa}
\usepackage{xfrac}
\usepackage[unicode,hidelinks,bookmarks=false]{hyperref}
\newcommand{\ie}{i.e.\ }
\newcommand{\cf}{cf.\ }
\newcommand{\resp}{respectively\ }
\newcommand{\Subsection}{\subsection}
\providecommand{\nobreakdash}{\nobreak\hbox{-}}
\setlength{\emergencystretch}{2em}
\multlinegap0pt
\usepackage{bm}
\usepackage{bbm}
\usepackage{dsfont}
\usepackage{xspace}
\usepackage{cancel}
\usepackage{mathtools}
\usepackage{amsthm}
\makeatletter
\@ifundefined{theorem}{\newtheorem{theorem}{Theorem}}{}
\makeatother
\DeclareMathOperator{\Gal}{Gal}
\newcommand{\FF}{{\mathbf{F}}}
\title[Locally imprimitive points on elliptic curves]{Locally imprimitive points on elliptic curves}
\author{Nathan Jones, Francesco Pappalardi, Peter Stevenhagen}
\date{Source: arXiv:2304.03964v1 (CC BY 4.0)}
\begin{document}
\maketitle

\noindent\emph{Source and licence.} Locally imprimitive points on elliptic curves. Version arXiv:2304.03964v1 (CC BY 4.0), distributed under the Creative Commons Attribution 4.0 International licence, as recorded in the arXiv licence metadata for this exact version and shown on the abstract page https://arxiv.org/abs/2304.03964v1. Full rights notice: licenses/arxiv-2304.03964-CC-BY-4.0.txt.\par\smallskip

\noindent\textit{Editorial scope.} Counted unit: the Theorem whose source label is \texttt{ell=3} in the article (source0.tex lines 1386--1402, source label \texttt{ell=3}), delivered together with the article's own complete original proof of it (source0.tex lines 1403--1454, one \texttt{proof} environment of 52 lines). No mathematics is written, repaired or re-derived: statement and proof are the article's own text line for line. Required context retained uncounted: 3-isogcurve-ab (lines 1361--1364); explicit3isog (lines 1367--1372); Deuringnf (lines 1379--1382). Every definition, display and label of those lines is retained unchanged. Omitted from this package and not redistributed: 1--1360, 1365--1366, 1373--1378, 1383--1385, 1455--2009. Nothing inside a retained excerpt is omitted or altered: every retained body is delivered exactly as the article's lines are written, so no omission receipt of any kind accompanies this package. Citations: the retained excerpts carry no citation command at all, so no bibliography entry of the article is transcribed and no reference is redistributed. Reference adaptation: none was needed and none was made; every reference inside the delivered unit resolves inside it, so no unresolved reference or citation is produced. Printed slips kept literal: no printed word, symbol, hypothesis or number of the counted statement or of its proof was altered. Layout and class: the article's own class is \texttt{amsart} with packages including amssymb, amsrefs, hyperref, booktabs, xy, multirow; the wrapper is the lane's pinned \texttt{amsart} editorial wrapper at 10pt on A4, which restates the theorem-like environments the retained text needs and adds the article's own macro definitions that the retained text needs (\texttt{\textbackslash FF}, \texttt{\textbackslash Gal}), with extra wrapper packages (bm, bbm, dsfont, xspace, cancel, mathtools). The delivered numbering is the unit's own. Assets: no retained span contains a figure environment, a graphics-inclusion command, a PDF-page inclusion, a table or a TikZ picture; no image file is copied into this package, the package declares no asset dependency and no image file is redistributed with it. Licence: version arXiv:2304.03964v1 of this article is distributed under the Creative Commons Attribution 4.0 International licence, as recorded in the arXiv licence metadata for this exact version and shown on the abstract page https://arxiv.org/abs/2304.03964v1. A full rights notice is delivered at licenses/arxiv-2304.03964-CC-BY-4.0.txt. Characters: every retained excerpt is pure ASCII, so the delivered engine, XeLaTeX with its default Latin Modern text font, reports no missing character.
\begin{equation}
\label{3-isogcurve-ab}
E: y^2+axy+by=x^3-5abx-(a^3+7b)b,
\end{equation}

\begin{equation}
\label{explicit3isog}
\varphi_3(x,y)= \left(
\frac{x^3+abx+b^2}{x^2},
    \frac{y(x^3-abx-2b^2)-b(ax+b)^2}{x^3}\right).
\end{equation}

\begin{equation}
\label{Deuringnf}
E_b': y^2+xy+by=x^3.
\end{equation}

\begin{theorem}
\label{ell=3}
For $K$ a number field, take $b=b(X)=(1-X-X^2)/X\in K(X)$ and 
define the associated elliptic curve over $K(X)$ as
$$
E_{b}: y^2+xy+by=x^3-5bx-(1+7b)b.
$$  
Then for infinitely many $t\in K^*$, the specialization $E_{b(t)}$
of $E_b$ is an elliptic curve over $K$ for which
\begin{equation}
\label{pointPt}
    \left(\frac{(t^2+t)(t^2-1)+1}{t^2},
    \frac{(1-t^2)\left(t^4+2 t^3+t-1\right)}{t^3}\right)
\in E_{b(t)}(K)
\end{equation} 
is a non-trivial locally 3-imprimitive point.
\end{theorem}

\begin{proof}
For the curve $E'_b$ in \eqref{Deuringnf} with $b\in K(X)$ as defined, 
the point $P'_b=(1, X)$ lies on $E'_b$ by the very choice of $b$:
it satisfies $X^2+X+bX=1$.
Under the 3-isogeny $\varphi_3: E'_b\to E_b$ from 
\eqref{explicit3isog} to the curve
$E_b$ obtained by putting $a=1$ in \eqref{3-isogcurve-ab},
it is mapped to 
$$
P_b=(1+b+b^2, (1-b-2b^2)X-b(1+b)^2)\in E_b(K(X)).
$$
Under the specialization $X=t\in K^*$, we obtain a point $P_{b(t)}$,
on the curve $E_{b(t)}$ defined over $K$ 
that is given by \eqref{pointPt}.
We are only interested in those 
specializations for which $E_{b(t)}$ is an elliptic curve.
As these are all $t\in K^*$ for which $b(t)\notin \{0, 1/27\}$, at most
4 `bad' values of $t$ are excluded.
Also, by the same argument as we gave for $P_\lambda$ in the case $\ell=2$,
there are only finitely many $t\in K^*$ for which $P_t$ is torsion point. 
These finitely many $t$ we also exclude as `bad' values.

We saw already that $E_b$ has 3-division field 
$K(\zeta_3, \root 3 \of b)$, and an explicit computation 
shows that the 3-division field of the point $P_b\in E_b(K(X))$ equals 
\[
K(X)(\textstyle\frac13 P_b)=K(\zeta_3, \root 3 \of b, \root 3 \of X)=
K(\zeta_3, \root 3 \of X, \root 3 \of {X^2+X-1}).
\]
Over $K(\zeta_3, X)$, the elements $X$ and $X^2+X-1$ have 
`independent' cube roots -- it suffices look at their ramification locus.
It follows that the Galois group of 
the 3-division field of $P_b$ over $K(X)$
may be described as
\[
\Gal(K(X)(\textstyle\frac13 P_b)/K(X))= 
\left\{ \begin{bmatrix}
\omega_3 & x & y \\
0 & 1 & 0 \\
0 & 0 & 1 
\end{bmatrix}: x,y\in\FF_3  \right\},
\]
with $\omega_3$ denoting the $\FF_3^*$-valued character corresponding to
the (possibly trivial) extension $K(X)\subset K(X, \zeta_3)$.
By Hilbert irreducibility, it follows that for infinitely many
$t\in K^*$ outside the finite set of `bad' values, 
the 3-division field of the point $P_{b(t)}\in E_{b(t)}(K)$ 
has the `same' Galois group over $K$,
making it into a point that is globally 3-primitive, but locally $3$-imprimitive.
As $E_{b(t)}$ does not have complete 3-torsion over $K$, we conclude that the point $P_{b(t)}$ given in \eqref{pointPt}
is a non-trivial locally 3-imprimitive point.
\end{proof}

\end{document}
